\section{A Sparse Set of \goalheads Carries the Goal}
\label{sec:heads}

\Cref{sec:what} read the goal out of a located set of attention heads. Here we say how those
\goalheads are found, and show that the goal signal is a property of the \emph{head selection}---not of
the residual stream at those positions, and not of a large fraction of the network.

\subsection{Locating the \goalheads}
\label{sec:heads:localize}

We localize the \goalheads with activation patching (\Cref{sec:notations:interventions}) on a
minimal-pair setup. A \emph{source} and a \emph{destination} prompt each ask the model to implement a
single predicate ($\pred_\src \neq \pred_\dst$) on the same data type (\texttt{int} $\mid$ \texttt{str}
$\mid \cdots$), each cut off so the next token is representative of the predicate being implemented. To
guarantee the located latents carry the abstract predicate rather than shallow token information, source
and destination use \emph{different programming languages}, so the source's target token is absent from
the destination's language (a concrete example is in {\color{red}\Cref{fig:goal_head_loc}}).

We cache what each head \emph{writes} to the residual stream (its OV contribution $a_t^{\ell j}$,
Eq.~\ref{eq:head}) from the source run, and learn a sparse binary mask over all heads that, when their
outputs are patched into the destination run, maximizes the logit of the source-predicate token
$\mathrm{tok}_\mathrm{targ}$:
\vspace{-1.0em}
\begin{align}
    \label{eqn:opt-int}
    a^{\ell j}_{-1} & \leftarrow \mathrm{mask}^{\ell j}\, a^{\ell j}_{\src} + (1 - \mathrm{mask}^{\ell j})\, a^{\ell j}_{-1}.
\end{align}
\vspace{-1.5em}

The learned mask selects a compact set of \goalheads concentrated in the model's middle-to-late layers
(band $25$--$56$ for \texttt{Qwen2.5-Coder-32B}; $122$ heads at threshold $>0.9$). Full training details,
layer/threshold sweeps, and per-model head sets are in \Cref{app:head_loc}.

\subsection{The head \emph{selection}, not the location, carries the goal}
\label{sec:heads:residual}

Is the goal simply present in the residual stream at the pivot, or specifically in what the \goalheads
write? We compare \goalscope's read-out quality when the injected vector is the \goalheads' aggregate
output versus the raw residual stream (and the full attention-block output) at the same positions and
depths, holding the injection mechanism fixed (norm-relative, so the comparison reads content, not
magnitude; \Cref{app:goalscope}).

\begin{table}[t]
\centering
\small
\begin{tabular}{lcc}
\toprule
\textbf{source injected into \goalscope} & \textbf{faithfulness} & \textbf{causality} \\
\midrule
\goalheads (122, output aggregate)            & \textbf{0.879} & \textbf{0.792} \\
full attention-block output, L40--50          & 0.343 & 0.334 \\
residual stream, L25--56                       & 0.031 & 0.020 \\
residual stream, L40--50                       & 0.009 & 0.006 \\
\bottomrule
\end{tabular}
\caption{The head selection does nearly all the work. \goalscope read-out on the locked
prefill-to-frontier corpus (Qwen2.5-Coder-32B, $n{=}253$, $\alpha{=}0.1$), varying only the injected
source. The goal is highly readable from the \goalheads' output, only partially from the full attention
output, and essentially \emph{unreadable} from the raw residual stream, whose norm is dominated by
non-goal content.}
\label{tab:heads-vs-residual}
\end{table}

\Cref{tab:heads-vs-residual} (see also \Cref{fig:heads-vs-residual}) shows the goal is highly readable
from the \goalheads (0.88 / 0.79) but essentially \emph{unreadable} from the raw residual at the same
depths ($\sim$0.01--0.03). The head selection, not the location, is what exposes the goal.

\begin{figure}[t]
\centering
\includegraphics[width=0.6\linewidth]{Figures/goalscope_axis_a_residual.png}
\caption{\goalheads{} aggregate vs.\ raw residual / attention-block output as the \goalscope source
(placeholder; regenerate at final quality).}
\label{fig:heads-vs-residual}
\end{figure}

\subsection{How many heads? The encoding is sparse}
\label{sec:heads:sparsity}

The read-out saturates with only a small fraction of heads (\Cref{tab:sparsity},
\Cref{fig:head-budget}). Faithfulness climbs steeply from $5$ to $25$ heads ($0.45 \to 0.79$) and then
plateaus: $\sim$25 heads ($\approx$0.6\% of the model's $4160$ heads) already recover $\sim$90\% of the
full-set read-out, and $50$--$100$ heads essentially match the $122$-head anchor. Adding heads beyond
$\sim$100 yields no further faithfulness and a mild \emph{causality} decline, consistent with larger
sets pulling in less-selective heads that dilute the steered direction. This sparsity is what makes
\goalscope and the steering experiments of \Cref{sec:steer} a targeted intervention rather than a
diffuse perturbation.

\begin{table}[t]
\centering
\small
\begin{tabular}{lcccccccccc}
\toprule
\# heads & 5 & 10 & 25 & 50 & 100 & \textbf{122} & 150 & 200 & 250 & 300 \\
\midrule
faithfulness & 0.446 & 0.715 & 0.791 & 0.817 & 0.867 & \textbf{0.879} & 0.859 & 0.880 & 0.852 & 0.874 \\
causality    & 0.425 & 0.625 & 0.779 & 0.839 & 0.813 & \textbf{0.792} & 0.816 & 0.794 & 0.750 & 0.779 \\
\bottomrule
\end{tabular}
\caption{\goalscope read-out vs.\ number of \goalheads (Qwen2.5-Coder-32B, \texttt{additive\_norm}
$\alpha{=}0.1$, \texttt{hard\_yes} pooled over the 5 pivots, $n{=}1265$ records). The read-out saturates
by $\sim$50--100 heads; the $122$-head set used elsewhere sits on the plateau.}
\label{tab:sparsity}
\end{table}

\begin{figure}[t]
\centering
\includegraphics[width=0.6\linewidth]{Figures/fig3_250_vs_122.png}
\caption{Read-out quality vs.\ number of \goalheads (placeholder). A handful of heads recovers most of
the signal.}
\label{fig:head-budget}
\end{figure}
